\documentclass[12pt,reqno]{amsart}
\usepackage{math-review}
\newcommand{\ii}{\sqrt{-1}}
\newcommand{\ddc}{\ii\partial\bar\partial}
\newcommand{\K}{\mathcal K}
\title[dHYM on surfaces]{The dHYM equation on surfaces:\\a positivity criterion and its mechanism}
\date{Compact reference adaptation --- 24 September 2026}
\newcommand{\SurfaceStatement}{%
Let $X$ be a compact K\"ahler surface, $\Omega\in\K_X$, and
$A\in H^{1,1}(X,\mathbb R)$ with $s=A\cdot\Omega>0$. Set
\[
 \widehat\theta=\operatorname{Arg}\int_X(\Omega+\ii A)^2\in(0,\pi),
 \qquad c=\cot\widehat\theta=\frac{\Omega^2-A^2}{2s}.
\]
The following are equivalent:
\begin{enumerate}
\item For every K\"ahler form $\omega\in\Omega$, there is a smooth closed real
$(1,1)$-form $\eta\in A$ such that
$\sum_{j=1}^2\arctan\lambda_j(\eta;\omega)=\widehat\theta$.
\item $B:=A+c\Omega$ is a K\"ahler class.
\item $B\cdot C>0$ for every irreducible curve $C\subset X$.
\end{enumerate}
}
\begin{document}
\begin{abstract}
On a compact K\"ahler surface, positive-phase scalar dHYM solvability is
an ordinary K\"ahler-cone condition after a cohomological shift. We state
the full criterion and derive its reduction to the complex Monge--Amp\`ere
equation, including the converse phase check. This is a surface-focused
expository sample, not a comprehensive literature update.
\end{abstract}
\maketitle
\section{Introduction}
Fix a K\"ahler form $\omega$ on a compact complex surface $X$.
For a smooth real $(1,1)$-form $\eta$, let $\lambda_1,\lambda_2$ be its
$\omega$-eigenvalues. With each arctangent in $(-\pi/2,\pi/2)$,
\[
 (\omega+\ii\eta)^2
 =\prod_{j=1}^2(1+\ii\lambda_j)\,\omega^2,
 \qquad \vartheta_\omega(\eta)=\sum_{j=1}^2\arctan\lambda_j.
\]
For a fixed closed representative $\alpha\in A$, the unknown is a real
smooth potential $\varphi$, with $\eta=\alpha+\ddc\varphi$.
The scalar dHYM equation is $\vartheta_\omega(\eta)=\widehat\theta$.
The positive phase fixes the lift, not merely the phase modulo $\pi$.
Here $D\cdot E=\int_XD\wedge E$ and $\K_X$ is the K\"ahler cone.

\begin{theorem}\label{thm:surface}
\SurfaceStatement
\end{theorem}
The analytic reduction is due to Jacob--Yau \cite{JY}; the last condition
uses the numerical description of the surface K\"ahler cone \cite{DP}.
The strict curve inequality is essential: nefness is not the asserted
solvability criterion.

\section{The shift and the Monge--Amp\`ere equation}
\begin{theoremrecall}{thm:surface}
\SurfaceStatement
\end{theoremrecall}
Write $\alpha_j=\arctan\lambda_j$ and $\beta=\eta+c\omega$.
If $\alpha_1+\alpha_2=\widehat\theta\in(0,\pi)$, then
\[
 \lambda_j+c
 =\frac{\cos(\widehat\theta-\alpha_j)}
 {\sin\widehat\theta\cos\alpha_j}>0.
\]
Indeed $\widehat\theta-\alpha_j$ is the other angle in
$(-\pi/2,\pi/2)$. Thus $\beta>0$. Expanding the imaginary part gives
\begin{align*}
 \operatorname{Im}\bigl(e^{-\ii\widehat\theta}
                (\omega+\ii\eta)^2\bigr)=0
 &\iff \eta^2+2c\,\omega\wedge\eta=\omega^2\\
 &\iff \beta^2=(1+c^2)\omega^2.
\end{align*}
In particular, a solution supplies a K\"ahler representative of $B$.

Conversely, suppose $B\in\K_X$ and choose $\beta_0\in B$ positive.
Cohomological compatibility is automatic:
\[
 B^2=A^2+2cs+c^2\Omega^2=(1+c^2)\Omega^2.
\]
Yau's theorem \cite{Yau} produces a smooth $\psi$, unique modulo constants,
with
\[
 \beta=\beta_0+\ddc\psi>0,
 \qquad \beta^2=(1+c^2)\omega^2.
\]
Let $\eta=\beta-c\omega$. Its class is $A$, and its eigenvalues satisfy
$(\lambda_1+c)(\lambda_2+c)=1+c^2$ with both factors positive. Hence
\[
 \lambda_1+\lambda_2
 \geq 2\sqrt{1+c^2}-2c>0.
\]
Therefore $\operatorname{Im}(\omega+\ii\eta)^2>0$.
Together with the vanishing imaginary part after rotation, this selects
$\vartheta_\omega(\eta)=\widehat\theta$, not another lift.
The $\partial\bar\partial$-lemma expresses $\eta$ in the prescribed
potential class. This proves the first two conditions equivalent.

The numerical step also needs the correct component of the positive cone.
By the Hodge index inequality, $A^2\Omega^2\leq s^2$, so
\[
 B\cdot\Omega
 =\frac{2s^2+\Omega^4-A^2\Omega^2}{2s}
 \geq\frac{s^2+\Omega^4}{2s}>0,
 \qquad B^2>0.
\]
On a compact K\"ahler surface, a real $(1,1)$-class with these two properties
is K\"ahler precisely when it pairs positively with every irreducible
curve \cite{DP}. The component condition has thus been checked rather
than suppressed.

\section{A model and the boundary of the comparison}
If $A=a\Omega$ with $a>0$, choose $\eta=a\omega$. Then
\[
 \widehat\theta=2\arctan a,\qquad
 c=\frac{1-a^2}{2a},\qquad
 B=\frac{1+a^2}{2a}\Omega\in\K_X.
\]
This model displays both the cohomological shift and its strict positivity.
For a general class, testing $A$ alone would miss the phase-dependent term.
Neither this calculation nor the existence theorem proves convergence of
an evolution starting at an arbitrary initial potential.

The mirror-motivated organizing question is the Thomas--Yau problem
\cite{TY}. For a Calabi--Yau metric with normalized parallel holomorphic volume form
$\Upsilon$, a graded Lagrangian $L$ has a real phase lift $\theta_L$ defined by
$\Upsilon|_L=e^{\ii\theta_L}\,dV_L$. A special Lagrangian has constant
$\theta_L$. The existence question asks which appropriate stability
conditions select Hamiltonian classes containing such a representative.
The flow question additionally asks what evolution and treatment of
singularities reach that representative. These are different assertions;
Joyce's categorical and surgery program \cite{Joy} is not a theorem of
unconditional smooth convergence.

The surface theorem above is complete in its stated scalar scope. It does
not identify general categorical stability with dHYM solvability, settle
Thomas--Yau, or extend the surface curve test to arbitrary dimension.
Those threads belong to the full reference article rather than to a
compressed statement with hidden hypotheses.

\begin{thebibliography}{JY17}
\bibitem[DP04]{DP} J.-P. Demailly and M. P\u{a}un, \emph{Numerical characterization of the K\"ahler cone of a compact K\"ahler manifold}, Ann. of Math. \textbf{159} (2004), 1247--1274.
\bibitem[JY17]{JY} A. Jacob and S.-T. Yau, \emph{A special Lagrangian type equation for holomorphic line bundles}, Math. Ann. \textbf{369} (2017), 869--898; \href{https://arxiv.org/abs/1411.7457}{arXiv:1411.7457}.
\bibitem[Joy15]{Joy} D. Joyce, \emph{Conjectures on Bridgeland stability for Fukaya categories of Calabi--Yau manifolds, special Lagrangians, and Lagrangian mean curvature flow}, EMS Surv. Math. Sci. \textbf{2} (2015), 1--62; \href{https://arxiv.org/abs/1401.4949}{arXiv:1401.4949}.
\bibitem[TY02]{TY} R. P. Thomas and S.-T. Yau, \emph{Special Lagrangians, stable bundles and mean curvature flow}, Comm. Anal. Geom. \textbf{10} (2002), 1075--1113; \href{https://arxiv.org/abs/math/0104197}{arXiv:math/0104197}.
\bibitem[Yau78]{Yau} S.-T. Yau, \emph{On the Ricci curvature of a compact K\"ahler manifold and the complex Monge--Amp\`ere equation, I}, Comm. Pure Appl. Math. \textbf{31} (1978), 339--411.
\end{thebibliography}
\end{document}
